\section{Binary testing and average correction at a fixed radius}
\label{sec:inner-binary}

We convert the polynomial arrays of Section~\ref{sec:inner-algebra} into
binary words.  The correction guarantee is averaged over the requested
encoded coordinates.  Its radius is fixed before the requested accuracy
and the density of those coordinates are specified.  This quantifier
order is essential when the one-port construction is subsequently used
many times.

All notation in this section is inner notation.  In particular, the
groups below are polynomial-size affine groups and are unrelated to the
outer index group.  Fix one of the padded templates of
Lemma~\ref{lem:inner-algebra-interface}.  We call each polynomial family
in its canonical tuple, together with its slot axes, a \emph{track}.
We use its fields
\(\F_q\subset K=\F_{q^\ell}\), common target domain \(K^u\), degree
bound \(D=\ell^2-1\), and fixed number of slot-mask axes per track.
The integer \(u\) and the field conventions are common to the actual
and one-port templates.  Honest chart degrees, including the constant
factor enlargement below, are \(O(\ell h)<D\) for sufficiently large
size.  The field-size choice already made in
Section~\ref{ia:path-parameters} ensures \(q/D\longrightarrow\infty\).
Constants in this section may depend on the fixed template, but are uniform
in the public circuit within that template.

\subsection{Rectangular coordinates and the common row code}
\label{ib:rectangular-code}

The algebraic neighbor formulas use invertible affine maps of the target
\(K^u\). We will add parity rows so that local row decoding can reduce
errors, and use a group action to move a requested value among columns
without changing it. The rectangular coordinates below make this action
transitive on point and ordered-line columns. Target affine maps commute
with it, so neighbor updates remain column permutations after row encoding.

Choose a constant integer \(r>2u\), \(r\ge3\).  A rectangular point
is \(I=(P,z)\), where \(P(x)=Ax+b\) maps \(K^r\) affinely to
\(K^u\) and \(z\in K^r\).  Retain points with \(\rank A=u\).
For an ordered \(\F_q\)-line
\[
 I_t=(P_0+tP_1,z_0+tz_1),\qquad t\in\F_q,
\]
retain it precisely when the matrix obtained by stacking \(A_0,A_1\)
has rank \(2u\) and \(z_1\ne0\).  Every point on such a line is
retained, since a dependence in \(A_0+tA_1\) would give a dependence
in the stacked matrix.  A sequential row-exposure bound shows that
the omitted point and line fractions are \(o(1)\) as \(\abs K\to
\infty\); for example rank failure for an \(s\)-by-\(r\) uniform
matrix is at most \(\sum_{j=0}^{s-1}\abs K^{j-r}\).

Write
\[
 U(K)=K^r\rtimes\mathrm{SL}_r(K),\qquad
 \mathcal H=U(K)\times U(K),\qquad
 (g_1,g_2)(P,z)=(P\circ g_1^{-1},g_2z).
\]
This is an affine action on the rectangular \(\F_q\)-coordinates,
and acts on ordered lines without changing their parameter.  There
is one orbit of retained points and one orbit of retained ordered
lines.  Indeed, right special-linear transformations are transitive
on full-rank matrices with \(u\) rows, and on those with \(2u\)
rows: first align by a general-linear transformation and then adjust
its determinant on an unused complementary dimension.  Translation
in the argument of \(P\) adjusts the one or two offsets because the
corresponding matrix has full row rank.  Independently, the second
affine factor is transitive on points and on ordered pairs of a base
point and a nonzero direction.  All retention conditions are
invariant.  In particular the point at each fixed parameter of a
uniform retained ordered line is uniform on retained points.

An invertible affine target map \(d:K^u\to K^u\) induces
\((P,z)\mapsto(d\circ P,z)\).  It commutes with \(\mathcal H\),
preserves both retained domains, and preserves line parameters.  A
target chart \(f\) is represented by \(f(P(z))\); this substitution
increases its coordinate degree by at most a factor two.

\begin{lemma}[Uniform affine-group expansion]
\label{ib:affine-expansion}
There is a fixed symmetric labelled multiset
\(u_1,\ldots,u_{d_e}\) in each \(\mathcal H\), with \(d_e\ge8\),
whose averaging operator has norm at most \(1/64\) on mean-zero
functions.  The labels, group operations, and complete group tables
are deterministically constructible in polynomial time in the inner
size.  The degree and gap are independent of the fields.
\end{lemma}

\begin{proof}
Here are the two primary group-theoretic inputs and their exact use.
Ershov--Jaikin-Zapirain~\cite[Theorem~1.1]{EJZ2010} prove property
\((T)\) for \(\mathrm{EL}_r(R)\) for every finitely generated
associative unital ring \(R\) and \(r\ge3\).  Their generating-set
bound in Section~6.1 applies to the elementary matrices with entries
in a fixed ring-generating set containing 1.  In particular, their
equations (6.1)--(6.2) give the positive Kazhdan constant
\[
 \kappa_r=\bigl[8(48\sqrt2+2\sqrt{3r})\bigr]^{-1}
\]
for
\(\{e_{ij}(1),e_{ij}(t):i\ne j\}\) in
\(\mathrm{SL}_r(\F_2[t])\); the equality with the elementary group
holds because \(\F_2[t]\) is Euclidean.  No restriction excluding
characteristic two occurs in this theorem.

Shalom~\cite[Theorem~3.4 and Corollary~3.5]{Shalom1999}, applied to the discrete
commutative ring \(R=\F_2[t]\), bounds the displacement of every
translation in \(R^2\rtimes\mathrm{SL}_2(R)\) by 484 times the
maximum displacement on the basis translations and
\(e_{12}(1),e_{12}(t),e_{21}(1),e_{21}(t)\).  Signs coincide in
characteristic two.  This is the homogeneous displacement form of
the corollary's tolerance statement, obtained by scaling and taking
limits.  It bounds the translation sector; it alone does not assert
expansion of the full affine group.

Set \(\Lambda=R^r\rtimes\mathrm{SL}_r(R)\) and let \(S\) consist
of all basis translations and the elementary matrices just listed
for all ordered coordinate pairs.  Embed Shalom's rank-two group on
each such pair.  Since a translation is a product of \(r\)
coordinate translations, every unitary representation satisfies
\[
 \sup_{a\in R^r}\norm{av-v}
       \le 484r\max_{s\in S}\norm{sv-v}.
 \tag{\ref{ib:affine-expansion}.1}
\]
Let \(V_N\) be the translation-invariant subspace.  Its orthogonal
projection commutes with \(\Lambda\), since the translations are
normal.  On \(V_N^\perp\), the minimum-norm point of the closed
convex translation-orbit hull is zero.  The displayed bound therefore
gives
\(\norm v\le484r\max_{s\in S}\norm{sv-v}\) on this subspace.
On \(V_N\), a representation without group-invariant vectors is a
representation of \(\mathrm{SL}_r(R)\) without invariant vectors,
so the preceding Kazhdan bound applies.  Orthogonally combining the
two bounds gives a positive Kazhdan constant \(\kappa\) for
\((\Lambda,S)\).  Explicitly, if \(\kappa_r\) is the quoted
 linear-group constant, one may take
\(\kappa=((484r)^2+\kappa_r^{-2})^{-1/2}\).

The set \(S\) generates \(\Lambda\): the third-index elementary
commutator multiplies entries, entries in a root subgroup add, and
conjugating basis translations gives arbitrary coordinate
translations.  Choose a field presentation
\(K=R/(p(t))\) by a monic irreducible polynomial.  Elementary
generation over a field shows that the quotient map surjects onto
\(U(K)\).  Pulling back its mean-zero regular representation
preserves the Kazhdan bound.  If \(d=\abs S\), the half-lazy walk
has Dirichlet form
\[
 \langle v,(1-M)v\rangle
   =\frac{1}{4d}\sum_{s\in S}\norm{sv-v}^2
   \ge\frac{\kappa^2}{4d}\norm v^2
\]
on mean-zero vectors, and nonnegative spectrum.  On the product
group, choose a factor uniformly and take this step there.  The
commuting tensor-factor operators give another uniform positive
gap and nonnegative spectrum.  A sufficiently large fixed power has
mean-zero norm at most \(1/64\).  Keep all walk labels, including
identities and coincident group elements, as a symmetric multiset;
repeat labels if necessary to ensure \(d_e\ge8\).  All matrices
have fixed dimensions over polynomial-size fields, so enumeration
and every stated computation are polynomial-time finite-array
operations.
\end{proof}

Form a bipartite graph with both vertex sets \(\mathcal H\), and
edges
\(\mathcal R=[d_e]\times\mathcal H\), where \((i,g)\) joins
left \(g\) to right \(u_i g\).  Its cross-part normalized second
singular norm is \(\lambda\le1/64\).  Put
\(c_e=\lfloor d_e/4\rfloor\).

\begin{lemma}[Systematic row extension]
\label{ib:row-extension}
There is a polynomial-size extension field \(J/\F_q\) and two
\(c_e\)-by-\(d_e\) matrices \(H_L,H_R\) over \(J\) whose every
\(c_e\)-column minor is nonzero, such that the edge code
\begin{align}
 \sum_i(H_L)_{bi}U_{i,g}&=0,\label{ib:left-row}\\
 \sum_i(H_R)_{bi}U_{i,u_i^{-1}x}&=0\label{ib:right-row}
\end{align}
has a unique linear extension \(E\) from labels \(i>2c_e\) to all
labels.  The same holds with symbols in any finite-dimensional
\(J\)-space.  Each local code has distance at least \(c_e+1\),
and \(E\) commutes with every right shift
\((R_hU)_{i,g}=U_{i,gh}\).
\end{lemma}

\begin{proof}
Treat the \(2c_ed_e\) matrix entries as indeterminates.  The square
system matrix on parity labels \(1,\ldots,2c_e\) has nonzero
determinant polynomial: specialize left row \(b\) to select label
\(b\), and right row \(b\) to select label \(c_e+b\).  This gives
a permutation matrix, including when some generators coincide.
Multiply this determinant by all maximal minors of both local
matrices.  Each factor is a nonzero polynomial; their product is
nonzero and has total degree at most
\[
 2c_e\abs{\mathcal H}
     +2c_e\binom{d_e}{c_e}.
\]
Take the least extension \(J/\F_q\) of cardinality exceeding this
bound.  It has polynomial cardinality.  A nonzero polynomial of
smaller total degree than \(\abs J\) cannot vanish everywhere on
\(J^{2c_ed_e}\).  Exhaustive search over this constant number of
entries, testing the determinant and minors, therefore finds a
specialization in polynomial time.  Inverting the parity matrix
constructs \(E\).

Column independence proves local distance \(c_e+1\), also for
vector symbols: a nonzero vector-valued local word supported on at
most \(c_e\) columns would contradict that independence after any
nonzero coordinate projection.  Right shifts permute left stars
and right stars in \eqref{ib:left-row}--\eqref{ib:right-row} and
preserve the systematic labels.  Uniqueness proves \(ER_h=R_hE\).
\end{proof}

Choose this field, these matrices, and their deterministic ordering
from the common point and field parameters alone, independently of
the number or identities of tracks.  Every track uses this one
encoder.  On systematic rows prescribe, slotwise,
\[
 Y_{i,g}(I)=f(P'(z')),\qquad (P',z')=gI,
\]
and extend columnwise by \(E\).  Equivariance yields
\begin{equation}
 Y_{i,gh}(I)=Y_{i,g}(hI)
 \quad\hbox{for all rows, including parity rows.}
 \label{ib:covariance}
\end{equation}
Every parity-slot function is a degree-in-bound polynomial over
\(J\).  At an ordered retained line store its actual restriction,
using coefficients through degree \(D\) and padding by zero.
These coefficients obey the same row constraints and covariance.
The commuting target maps act by the same column permutations on
point tables, line tables, and their encoded parity rows.
In particular, covariance gives
\(Y_{i,gh}(h^{-1}I)=Y_{i,g}(I)\) for a point or an ordered line \(I\).
Thus changing the row and making the inverse change to the column
preserves the requested symbol. This is the movement among columns that
will accompany local row decoding.

\subsection{Binary blocks and smooth probes}
\label{ib:binary-blocks}

For each track there are point and line blocks, both denoted \(Y\)
with their type understood.  A point position specifies a row edge,
a retained point, one \(\F_q\)-vector mask per slot axis, and an
additive binary form \(\mu:J\to\F_2\).  Its bit is \(\mu\) of
the slot contraction.  A line position additionally specifies masks
\(A,B\in\F_q^\ell\) for the coefficient index.  If \(Y(t)\)
is already contracted on the original slots, the bit is
\begin{equation}
 \mu\left(\sum_{j,l=0}^{\ell-1}A[j]B[l]
                     [t^{j+\ell l}]Y(t)\right).
 \label{ib:line-bit}
\end{equation}
This uses \(D+1=\ell^2\).  Evaluation at a specified parameter
uses the separable masks \((t^j)_j\) and \((t^{\ell l})_l\),
with the usual convention \(t^0=1\).

The short two-mask coefficient representation in \eqref{ib:line-bit}
is adapted, through the inherited polynomial-size construction, from
Amireddy--Behera--Srinivasan--Sudan--Willumsgaard
\cite[Section~3.1, Definition~3.1 and Claim~3.2]{ABSsw26}.
This attribution concerns the coefficient encoding and its evaluation
masks, not the correction and evolution machinery developed here.

Fix an \(\F_q\)-linear map \(\pi:J\to\F_q\) restricting to
the identity on \(\F_q\).  For each compatible ordered pair of
point-track types, include a product block \(Z\).  Compatibility
means that one tag is zero or both positive tags agree.  A type may
be paired with itself, with two independent positions.  An original
\(Z\) position consists of the two complete edge/point/slot-mask
positions, \(b_1,b_2\in J\), and a binary bilinear form \(Q\)
on the \(\F_2\)-coordinates of \(\F_q\).  For the two contracted
symbols the bit is
\begin{equation}
 Q\bigl(\pi(b_1Y_1),\pi(b_2Y_2)\bigr).
 \label{ib:product-bit}
\end{equation}
There are no line-product blocks and no products of distinct
positive tags.

Give every block equal weight.  If its original length is \(L_b\),
repeat each original bit uniformly until every block has length
\(\prod_b L_b\).  Every raw query chooses a uniform repetition
copy.  Let \(b_\#\) be the block count, \(m_{\rm in}\) the total
length, and \(\mathcal W(v)\) the canonical word for an extended
label \(v\).  These lengths are polynomial in the inner size:
there are constantly many blocks, \(q^{O(\ell)}\) slot masks,
fixed-dimensional matrix domains over \(K\), polynomial-size
\(J\), and \(2^{(\log_2q)^2}\) bilinear forms \(Q\).
Here \(\log q=O(\log\log n')\), so the last quantity is also
polynomial in \(n'\).  The product used for balancing has constantly
many polynomial factors.  Lengths are fixed within a padded template
independently of the circuit coefficients.  Positive-tag original
indices and canonical formulas coincide with the corresponding
one-port indices; repetition counts need not coincide.

We will repeatedly use the following explicit smoothing operation.
For a specified additive argument \(a\), choose uniform \(r\)
and write \(a=r+(a+r)\).  Multilinearity expands a requested value
into a bounded XOR of terms, each uniform in that argument.
Split different axes independently.  This applies to slot masks,
coefficient masks, bit forms, \(Q\), and, when needed, the two
multiplier axes.  Each term has uniform marginals jointly across
its independently split axes, conditional on the entire specified
request.  Pull back all public scalars and other linear maps before
splitting.  This remains true for zero coefficients; a zero term
may instead be omitted and supplied as public zero.

To probe a target value at \(y\in K^u\), choose a uniform row
group \(g\), full-rank \(A\), and \(z\), set
\(P(x)=Ax+y-Az\), and use the systematic position
\begin{equation}
 (i,g,g^{-1}(P,z)).
 \label{ib:systematic-probe}
\end{equation}
For uniform \(y\), this is jointly uniform in row and retained
point within its systematic label; every fiber has the same number
of \((A,z)\) choices.  A density bound for \(y\) therefore gives
the same bound for this joint distribution.  Independently chosen
target arguments can use independent representations.  An additive
form on a target \(\F_q\)-value extends to \(J\), for example by
composition with \(\pi\).  These procedures are defined on every
raw word and are smooth before any replacement or soundness claim.

\subsection{Dimension-independent replacement lemmas}
\label{ib:replacement-tools}

The tester's soundness proof replaces arbitrary binary data by symbol
tables and then by polynomial arrays. The two replacement lemmas below
make the error cost independent of the number of field coordinates and
slot entries.

We state precisely the low-degree input being used.  The corrected
Theorem~13 of Friedl--Sudan~\cite{FS1995} says that for every
\(\eta>0\) there is \(C_\eta<\infty\) such that, if
\(\abs{\F}\ge C_\eta D\), \(D\ge1\), and a function
\(f:\F^m\to\F\) and degree-at-most-\(D\) univariate claims
\(p_{x,h}\) satisfy
\[
 e=\Pr_{x,h\in\F^m,\,t\in\F}
       [p_{x,h}(t)\ne f(x+th)]\le \tfrac18-\eta,
\]
then a polynomial \(p:\F^m\to\F\) of total degree at most \(D\)
satisfies \(\dist(f,p)\le2e\).  The variables in this probability
are independent and uniform, including the possible direction \(h=0\);
the field-size constant does not depend on \(m\).  We use
\(\eta=1/16\).  Neither a vector-valued conclusion nor a restriction
on the line domain is part of this cited statement.

\begin{lemma}[Vector and deleted-domain low-degree replacement]
\label{ib:vector-ld}
Let \(V\) be any finite-dimensional \(\F_q\)-space.  In the preceding
field regime, line--point disagreement tending to zero for
\(V\)-valued point functions and degree-in-bound line claims implies
point distance tending to zero from a single \(V\)-valued polynomial
of total degree at most \(D\), uniformly in both dimensions.
The conclusion remains true if the point and ordered-line domains
omit fractions tending to zero, every retained line lies entirely in
the retained point domain, and the point at a uniform parameter on a
uniform retained line is uniform on retained points.  In this case the
mean fraction of line coefficient vectors differing from the
polynomial's restrictions also tends to zero.
\end{lemma}

\begin{proof}
First use the full domains.  For every \(\alpha\in V^*\), scalar
projection and the cited theorem give a polynomial \(p_\alpha\) with
\(\dist(\alpha f,p_\alpha)\le2e\).  Distinct scalar polynomials in
this degree range differ at a fraction at least \(1-D/q\) of points.
Consequently, when \(6e<1-D/q\), comparison with
\((\alpha+\beta)f\) forces
\[
 p_{\alpha+\beta}=p_\alpha+p_\beta,
 \qquad p_{a\alpha}=a p_\alpha \quad(a\in\F_q).
\]
The second identity uses at most \(4e\) disagreement; the assertions
with a zero form or scalar follow in the same way.  Coefficientwise
finite-dimensional duality now gives one \(V\)-valued polynomial
\(p\) for which \(\alpha p=p_\alpha\) for all \(\alpha\).
A uniform \(\alpha\) detects each nonzero vector with probability
\(1-1/q\), so
\[
 (1-1/q)\dist(f,p)
 =\E_\alpha\dist(\alpha f,p_\alpha)\le2e.
\]
There is no factor depending on \(\dim V\).

For deleted domains, extend the point function arbitrarily and give
every added line any degree-in-bound claim, for example zero.  The
full incidence disagreement is at most the retained-line disagreement
plus the omitted-line fraction.  The full-domain conclusion therefore
gives vanishing error also on retained points.  On any retained line
whose claim differs from the restriction of \(p\), a nonzero scalar
projection of their difference is a univariate polynomial of degree
at most \(D\).  The two claims disagree at a fraction at least
\(1-D/q\) of its parameters.  Averaging the original test error and
the point replacement error, using the asserted uniform point
marginal, bounds the fraction of such lines by their sum divided by
\(1-D/q\).  This tends to zero and proves the coefficient claim.
\end{proof}

The addition-test argument below uses the majority self-correction method
of Blum, Luby, and Rubinfeld~\cite[Generic Linear Self-Testing]{BLR1993}.
We include it together with the field-homogeneity, multiple-axis, and
slice-averaging extensions needed for our vector alphabets.

\begin{lemma}[Axis replacement]
\label{ib:axis-replacement}
For elementary binary vector groups, the addition test has a
dimension-independent high-agreement replacement by an additive map,
also with a vector-group alphabet.  For \(\F_q\)-spaces of
characteristic two, addition and scalar-homogeneity tests give
replacement by an \(\F_q\)-linear map.  For any fixed number of axes,
testing these identities in each axis at uniform other arguments gives
replacement by a multilinear map as the failure probability tends to
zero, uniformly in every domain and alphabet dimension.  These
conclusions also hold in mean over additional slice parameters.
\end{lemma}

\begin{proof}
Suppose first that \(f:G\to A\) fails addition with probability
\(e\), where both groups have exponent two.  For fixed \(x\) put
\(D_x(y)=f(x+y)+f(y)\).  The addition identities on the pairs
\((x+y,y+z)\) and \((y,y+z)\) imply \(D_x(y)=D_x(z)\).
Each pair is uniform when \(y,z\) are independent uniform elements,
so the collision probability is at least \(1-2e\), for every \(x\).
The most frequent value \(L(x)\) therefore has mass at least
\(1-2e\).  The identity
\[
 D_{x+y}(z)=D_x(y+z)+D_y(z)
\]
and three majority bounds imply \(L(x+y)=L(x)+L(y)\) when
\(6e<1\): at least one \(z\) realizes all three majority values.
If \(f(x)\ne L(x)\), the conditional addition failure at \(x\)
is at least \(1-2e\).  Thus
\(\dist(f,L)\le e/(1-2e)\).

For the field version, transfer the homogeneity test to \(L\).
The scalars \(a\) satisfying \(L(ax)=aL(x)\) identically form an
additive subgroup of \(\F_q\).  If it is proper, at least half the
nonzero scalars fail; for each failing scalar its nonzero additive
discrepancy is nonzero on at least half the inputs.  The homogeneity
test with uniform nonzero scalar would then have failure at least
\(1/4\).  This contradicts failure tending to zero after a vanishing
replacement cost.  Hence \(L\) is field-linear.

Induct on the number \(j\) of axes.  Markov's inequality, with a
threshold tending slowly to zero, makes all but a vanishing fraction
of first-axis slices good for the remaining \(j-1\) tests.  Replace
good slices by multilinear maps and the remaining slices by zero.
This changes a vanishing fraction of values.  Regard the resulting
slices as a map into the vector space of multilinear maps on
\(j-1\) axes.  A nonzero such map is nonzero on at least
\(2^{-(j-1)}\) of uniform inputs: successively contract each axis,
using that a nonzero linear map has a kernel of relative size at most
\(1/2\).  Thus vanishing first-axis test failure on values gives
vanishing failure as equalities in this vector-space alphabet.
The one-axis result finishes the induction.  The same Markov cutoff
over any additional parameters proves the mean version.
\end{proof}

\subsection{Testing the binary representation}
\label{ib:testing}

\begin{proposition}[Smooth tester and distance]
\label{ib:tester}
For each fixed template, the canonical encodings of valid extended
labels have constant relative distance.  They admit a perfectly
complete bounded-bit tester with the following uniform gap: for every
fixed \(a>0\), there are \(\theta(a)>0\) and a size cutoff such that
rejection below \(\theta(a)\) implies distance less than \(a\) from
one canonical valid word.  Each potential raw read on each fixed
answer branch has bounded density relative to the balanced word,
without conditioning on execution of that branch.  Consequently
rejection is at most a fixed constant times distance from a canonical
valid word.
\end{proposition}

\begin{proof}
Take a positive fixed mixture of the test families described below,
including each present block type.  Their query counts are bounded,
and all requested nonuniform masks are implemented by the splitting
operation above.  We prove the gap by a sequential replacement
argument.  Consider any sequence of increasing sizes and words whose
rejection tends to zero, allowing the public circuits to vary in the
fixed template.  We show distance tending to zero from a canonical
valid word.  A bounded smooth test retains vanishing rejection after
a vanishing-distance change, by the union bound over its potential
reads; for a branching test use its entire bounded answer tree.

\paragraph{Copies, symbols, and slot tensors.}
Compare two independent random copies of a uniform original bit.
If the fraction of ones among its copies is \(p\), the test rejects
with probability \(2p(1-p)\), which is at least
\(\min(p,1-p)\).  Replacing all copies by their majority costs
vanishing distance.  In each point or line \(Y\) block test
additivity in \(\mu\).  Lemma~\ref{ib:axis-replacement} and binary
double duality replace the block, in mean over other indices, by
Hadamard tables of genuine \(J\)-symbols.  Next test addition and
\(\F_q\)-homogeneity in every slot-mask axis, including the two
coefficient-mask axes for a line.  Test symbol equality by uniform
\(\mu\), which detects any inequality with probability \(1/2\);
pull field scalars onto this form before smoothing its raw accesses.
The same replacement lemma makes all these symbols honest
multilinear slot tensors at each edge and point or line, at vanishing
mean cost.  No global polynomial property has yet been asserted.

\paragraph{Point and line polynomials.}
At a uniform edge, retained ordered line, and parameter, compare the
point symbol with line evaluation, using uniform original slot masks
and a uniform bit comparison.  Split the two specified evaluation
masks in \eqref{ib:line-bit}.  A nonzero symbol discrepancy has a
constant binary detection probability, so the line--point
disagreement tends to zero in mean.  For most edge and original-mask
tuples apply Lemma~\ref{ib:vector-ld}, with \(V=J\) as an
\(\F_q\)-space, to obtain a degree-at-most-\(D\) rectangular
polynomial.  Use zero on the exceptional tuples.  The lemma's
line-coefficient conclusion also replaces the line tables by the
restrictions of these polynomials at vanishing cost.

The selected polynomials need not yet be coherent between original
mask tuples.  Transfer the previous multilinearity identities to
them using their mean point error.  A false identity between
degree-in-bound polynomials is nonzero on at least
\(1-D/q-o(1)\) of retained points.  Thus the identities fail with
vanishing probability as equalities in the vector space of
degree-in-bound polynomials.  Apply
Lemma~\ref{ib:axis-replacement} in that vector-space alphabet, in
mean over edges, and use zero on exceptional edges.  Changing a
vanishing fraction of mask tuples changes a vanishing fraction of
point and line bits.  We now have genuine polynomial slot families
at every row, and their exact line restrictions.

\paragraph{Covariance and common systematic data.}
Test \eqref{ib:covariance} on point symbols with independent uniform
\(g,h,I\), uniform masks, and a bit form.  Undo each row's chart:
\[
 F_{i,g}(I)=Y_{i,g}(g^{-1}I).
\]
If two such polynomial slot tuples differ, the discrepancy is seen
with constant probability, first on a retained point, then by
successive nonzero mask contractions, then by the bit form.  Since
\(g\) and \(gh\) are independent uniform elements, for each fixed
label and track the squared frequencies of the uncharted tuples
sum to \(1-o(1)\).  Their maximum frequency is therefore
\(1-o(1)\).  Replace other rows by this predominant tuple, charted
back.  This changes vanishing point and line distance and enforces
covariance exactly, including on lines.

Compare uncharted tuples of the finitely many systematic labels,
using a shared uniform retained point and masks and independent row
groups.  The same constant gap forces an exactly common polynomial
tuple \(F\) on all systematic labels of a track.  Test that its
contracted values lie in \(\F_q\) by a random binary form on
\(J/\F_q\), pulled back to \(J\) and split.  The quotient is an
\(\F_q\)-space.  If a polynomial coefficient were outside
\(\F_q\), its quotient would be a nonzero vector polynomial and
would have a constant detection gap.  Consequently every systematic
coefficient is in \(\F_q\).

\paragraph{Factoring through the target and enforcing parity.}
Choose uniform \(y\in K^u\) and two conditionally independent
generic representations as in \eqref{ib:systematic-probe}.  Compare
their values, with independent row groups.  Each individual marginal
is smooth; this test is a linear comparison and makes no product
query on the correlated pair.  For a slot coefficient the uncharted
identity tested is
\begin{equation}
 F(A,y-Az,z)=F(A',y-A'z',z').
 \label{ib:factor-target}
\end{equation}
The difference has coordinate degree at most \(2D\), and the
full-rank restrictions omit an \(o(1)\) fraction.  A failed
polynomial identity therefore has constant detection probability,
also after the bounded slot contractions.  Hence
\eqref{ib:factor-target} holds identically.  Fix any full-rank
\(A'\) and any \(z'\).  Its right side defines a target polynomial
\(f(y)\) of degree at most \(D\) with \(\F_q\) coefficients.
Substitute \(y=Az+b\) to obtain the exact factorization
\(F(P,z)=f(P(z))\).

Test the left and right star equations at uniform stars and retained
points, with slot masks and random \(J\)-bit forms; pull row
coefficients onto these forms.  In an exactly covariant family, a
failed polynomial star relation fails on every right-shift translate
of that star, with an invertible affine change of argument.
Polynomial, mask, and bit detection thus give a constant rejection
gap for any failure.  All star identities hold once the rejection is
sufficiently small.  Lemma~\ref{ib:row-extension} now forces the
exact parity extension of the systematic target data and its line
restrictions.

\paragraph{Product blocks.}
In each present \(Z\) block test additivity in \(Q\) and replace
at vanishing cost by exact \(Q\)-linear maps, in mean over all its
other indices.  Sample a uniform independent pair of
edge/point/slot-mask indices and uniform multipliers.  Test the
product identity for
\(Q(x,x')=\alpha(x)\alpha'(x')\), with independent uniform
binary forms \(\alpha,\alpha'\) on \(\F_q\).
On the \(Z\) side split only \(Q\), keeping these already uniform
other indices fixed.  On each \(Y\) side pull the multiplier and
\(\pi\) into the requested bit form and split that form.  Compare
with the Boolean product of the two resulting bits.

The \(Y\)'s are already genuine symbols.  The difference between
a wrong \(Q\)-linear map and their true product is a nonzero binary
matrix.  A uniform outer-product form detects it with probability
at least \(1/4\), by two successive nonzero contractions.  The
fraction of wrong maps therefore tends to zero.  Replace them by
\eqref{ib:product-bit}.  This works separately for every compatible
pair and never invokes a missing product block.

\paragraph{Algebraic validity.}
Finally run the algebraic validator of
Lemma~\ref{lem:inner-algebra-interface} on the recovered target
polynomials, using systematic representations.  Linear target
probes need only their individual bounded-density marginal.  For
each product probe that lemma gives compatible tags and independent
uniform full target arguments.  Represent these independently in
rectangular coordinates.  The corresponding \(Z\) query therefore
has the required joint density.  Identity multipliers on
\(\F_q\)-symbols and the demanded bilinear form implement the
field-level product; split the necessary slot, form, and multiplier
axes for its raw reads.  The exact algebraic gap forces canonical
extended validity.  The native linear maps used by the validator,
including coordinate packing, need not commute with the row action:
they are used here only to specify systematic target arguments.

All replacements have vanishing total binary cost, so the original
word is at vanishing distance from a canonical valid word.  If the
stated fixed-gap conclusion failed, for each integer \(j\) one
could choose a size at least \(j\) and a word of distance at least
\(a\) with rejection below \(1/j\), contradicting this sequence
argument.  Each constituent test is perfectly complete for every
random choice in its support.  All its raw marginals are smooth by
the displayed representations and splitting; this also proves the
claimed linear upper bound on rejection near a canonical word.

For distance, different logical tapes differ in an order-one grid
polynomial.  Its difference is nonzero on a constant fraction of
uniform target arguments, hence on a constant fraction of systematic
row/point positions; a bit form detects the symbol difference with
probability \(1/2\).  A concerned block has fixed positive weight.
If the logical tape agrees, different valid extended labels differ
in mode or in the progress marker, because legacy and mode fix the
frame.  These constant slot tensors likewise have a fixed
nonzero-contraction gap.  This gives a constant lower bound on
relative distance.
\end{proof}

\subsection{Value-independent uncertainty supports}
\label{ib:uncertainty}

The correction argument uses uncertainty sets rather than the actual
wrong symbols at intermediate stages.  This permits separate
consistent choices in each bounded dependency cone.

The alternating local decoding strategy is classical; see
Z\'emor~\cite[Section~III-A, Lemma~5 and Theorem~6]{Zemor2001}.
The proofs below supply the vector-space, repeated-edge,
value-independent-support, and finite-cone refinements needed here,
rather than invoking that result verbatim.

For a \(Y\) block, a column fixes a retained point or line and all slot
masks, including line coefficient masks. At that column, varying the
edge gives a row codeword of \(J\)-symbols in the reference encoding.
The binary data for one such symbol form its Hadamard table, including
all repetition copies; we call this its \emph{inner table}. The first
lemma controls uncertainty in these row-code positions without examining
the uncertain symbols' values.

\begin{lemma}[One-axis support contraction]
\label{ib:one-axis-support}
For the row code above, set
\[
 t_* =\lfloor c_e/2\rfloor+1,\qquad
 \vartheta=t_*/d_e\ge1/8,\qquad \zeta_*=1/256.
\]
Starting with an uncertain edge set of density
\(\zeta\le\zeta_*\), a constant number of alternating local
decoding half-rounds leaves an arbitrarily small fixed uncertain
edge density.  The resulting support depends only on the starting
uncertainty set and contains every possible output error for every
assignment equal to the reference codeword off that set.  This holds
for symbols in any finite-dimensional \(J\)-space.
\end{lemma}

\begin{proof}
At a star decode to the unique local codeword within
\(t_*-1=\lfloor c_e/2\rfloor\) errors, if it exists; otherwise
output zero on that star.  Distance \(c_e+1\) gives uniqueness.
This decoder uses a bounded number of linear symbol operations and
symbol-equality decisions: enumerate the information sets of size
\(d_e-c_e\), solve using the complementary invertible minor,
and count agreements with each candidate.  There are constantly
many such choices.

On a half-round mark a star bad if it meets at least \(t_*\)
currently uncertain edges, and declare all edges incident to bad
stars possibly uncertain afterward.  Every other star decodes
correctly, regardless of the values on uncertain edges.  The first
bad-star fraction is at most \(\zeta/\vartheta\le8\zeta\).
If \(x\) and \(y\) are consecutive bad-star fractions on opposite
sides, the new bad stars receive at least a \(\vartheta\) fraction
of their incident edges from the preceding bad stars.  Expander
mixing with multiplicities gives
\begin{equation}
 \vartheta y\le xy+\lambda\sqrt{xy}.
 \label{ib:mixing-contraction}
\end{equation}
When \(x\le\vartheta/2\), either \(y=0\), or division and
squaring imply
\[
 y\le(2\lambda/\vartheta)^2x\le x/16.
\]
Since \(8\zeta_*\le\vartheta/2\), induction applies at every
subsequent half-round.  The current edge-support fraction equals
the latest bad-star fraction.  A constant number of half-rounds
therefore gives any prescribed fixed residual density.  Every
support operation used only sets of uncertain positions, proving
the value-independent assertion.
\end{proof}

For a \(Z\) block, a column fixes both retained points, their slot
masks, and \(Q\). Its row position is a joint pair of edges, and its
inner table consists of the multiplier pairs and repetition copies.
The reference multiplier tables obey row constraints along each edge
axis, but the two axes use different scalar actions. We verify these
actions before applying a two-axis contraction.

For \(Z\), the relevant symbol alphabet is the additive group
\(\mathcal B_{\rm mul}\) of all binary bilinear functions
\(F:J\times J\to\F_2\).  It carries two commuting \(J\)-space
structures
\begin{equation}
 (a\cdot_1F)(b_1,b_2)=F(ab_1,b_2),\qquad
 (a\cdot_2F)(b_1,b_2)=F(b_1,ab_2).
 \label{ib:two-actions}
\end{equation}
Additivity in each argument and the field laws verify the vector
space axioms and commutation.  Distinct such functions differ on
at least \(1/4\) of uniform multiplier pairs.  Fixing the two
outer points, the original slot masks, and \(Q\) makes the true
multiplier tables across edge pairs a two-axis codeword.  For
example the first-axis relation with coefficients \(a_j\) is
\[
 \sum_j Q\bigl(\pi(b_1a_jY_{1,j}),\pi(b_2Y_2)\bigr)
 =Q\bigl(\pi(b_1\sum_j a_jY_{1,j}),\pi(b_2Y_2)\bigr)=0.
\]
The second-axis relation is identical with the roles exchanged.
The next lemma permits precisely these two possibly different space
structures.

\begin{lemma}[Two-axis support contraction and cone consistency]
\label{ib:two-axis-support}
Consider an array on two edge axes whose reference obeys the row
code in each slice, with possibly different \(J\)-space structures
on the same additive symbol group for the two axes.  Put
\(\zeta_{**}=\zeta_*^2/4\).  From initial uncertainty density at
most \(\zeta_{**}\), a constant number of first-axis half-rounds
followed by a constant number of second-axis half-rounds leaves any
prescribed fixed residual density.  The uncertainty support is
value-independent.  For one output position, the same guarantee
holds when arbitrary values are chosen on initially uncertain
positions specifically for its finite dependency cone, provided
each repeated initial position receives one consistent value inside
that cone.
\end{lemma}

\begin{proof}
First-axis slices with uncertainty density greater than
\(\zeta_*\) occupy at most a \(\zeta_*/4\) fraction.  Keep these
whole stripes uncertain.  On all other slices use
Lemma~\ref{ib:one-axis-support} to make residual density at most
\(\delta_1\).  In all but \(2\delta_1/\zeta_*\) of the
second-axis slices, the additional residual outside these stripes
has density at most \(\zeta_*/2\).  Such a slice starts with
uncertainty at most \(3\zeta_*/4<\zeta_*\).  Contract it to
\(\delta_2\), leaving all other second-axis slices uncertain.
The final density is at most
\[
 2\delta_1/\zeta_*+\delta_2,
\]
which is arbitrarily small.  Intermediate arrays need not satisfy
both slice constraints: each phase is compared to the fixed
reference and uses its appropriate space structure only.

A fixed number of rounds and constant degree give a finite
dependency cone.  Extend any consistent assignment in that cone
arbitrarily to the other uncertain initial positions, retaining the
reference elsewhere.  The global support statement applies to this
extension and determines the same output.  Separate cones may use
different values on their uncertain inputs; no common assignment
between cones is needed.
\end{proof}

\subsection{Bit simulation of a finite decoding cone}
\label{ib:cone-simulation}

\begin{lemma}[Inner-table flags and arbitrary-form prediction]
\label{ib:inner-table-prediction}
There are constants \(0<\nu_1<\nu_2<1/64\), independent of the
number of row rounds, with the following properties.  For a table
of bits, including its actual repetition copies, one can distinguish
distance at most \(\nu_1\) from a legal inner table from distance
greater than \(\nu_2\) from every legal inner table, with any
prescribed fixed reliability and bounded smooth queries.  The legal
inner tables are either Hadamard tables of \(J\)-symbols or binary
bilinear functions on \(J\times J\).  Inside radius \(\nu_2\)
the legal symbol is unique, and any specified bit of it can be
predicted with arbitrarily high fixed reliability by bounded queries
with uniform inner-index and copy marginals.  The specified form or
multiplier pair may be any function of previously exposed data;
the assertion holds conditional on those data, using fresh randomness.
\end{lemma}

\begin{proof}
Random-copy agreement and binary addition, or the two-axis binary
tests, are perfectly complete.  Lemma~\ref{ib:axis-replacement},
together with the majority replacement of copies, gives a
dimension- and copy-count-independent positive rejection lower
bound \(\gamma_2\) at distance greater than a fixed \(\nu_2\)
from every legal inner table.  If this failed, a sequence of such
tables with rejection tending to zero would contradict that lemma.
Near a legal table, smoothness bounds rejection by \(C_2\) times
distance, with an absolute \(C_2\).  Choose
\(\nu_1<\min(\nu_2,\gamma_2/(4C_2))\).  Estimating rejection
from independent trials and comparing with \(\gamma_2/2\)
distinguishes the two promised cases with any fixed reliability;
Chebyshev's inequality suffices to choose a bounded sample count.
The intermediate annulus may be assigned either flag.

Distinct Hadamard tables differ on half their forms.  Distinct
bilinear functions differ on at least a quarter of multiplier
pairs.  Repetition preserves these distances, proving uniqueness
within \(\nu_2\).  For a specified Hadamard form \(\mu\), choose
uniform \(\alpha\), query random copies at \(\alpha\) and
\(\mu+\alpha\), and XOR.  Each query is uniform in its inner
index, so one trial errs with probability at most \(2\nu_2\).
For a bilinear table at specified \((b_1,b_2)\), choose independent
uniform \(r_1,r_2\) and use
\begin{align*}
 F(b_1,b_2)={}&F(r_1,r_2)+F(b_1+r_1,r_2)\\
             &+F(r_1,b_2+r_2)+F(b_1+r_1,b_2+r_2).
\end{align*}
Each pair is jointly uniform, with a fresh random copy, and the
error probability is at most \(4\nu_2<1/16\).  Independent
repetition and majority give any desired fixed error.  These bounds
hold for every specified bit, so remain valid after conditioning
on all prior data if the splits and copies are fresh.
\end{proof}

\begin{lemma}[Consistent bit simulation]
\label{ib:bit-simulation}
Fix a finite cone of one- or two-axis row decoding, a reference satisfying
the corresponding row constraints, and a requested bit of the output
symbol. For analysis, declare an initial position uncertain when its
received inner table is more than \(\nu_1\) from its reference table,
and propagate this set by the support rules above. With bounded bit
queries the requested bit can be simulated to any fixed error probability
\(\sigma>0\), on a consistent assignment of initially uncertain
symbols that agrees with the reference at every initial table
within \(\nu_1\) of its true inner table. If the output position lies
outside the value-independent support, the predicted bit equals its
reference bit except with probability \(\sigma\).  All hidden operations are bit equalities
and requested bits; field coefficients and address computations
are public.  The routine is bounded and defined on arbitrary words.
\end{lemma}

\begin{proof}
List the cone's initial positions and deduplicate repeated positions
within that cone.  Flag each distinct table once by
Lemma~\ref{ib:inner-table-prediction}, calling the far outcome
flagged, and assign flagged positions the zero symbol.  At an
unflagged table within \(\nu_2\) of a legal table, assign its
unique nearby symbol.  Choose the flag reliability so that, except
with a prescribed small probability, every table within \(\nu_1\)
of its true symbol is unflagged and no table farther than \(\nu_2\)
from all legal symbols is unflagged.  On this event the assignment
is well-defined and consistent inside the cone.  Certain inputs
receive their true symbols: a competing symbol within \(\nu_2\)
would contradict inner distance.  The remaining assignments,
including arbitrary choices in the gap annulus, are permitted by
Lemmas~\ref{ib:one-axis-support} and~\ref{ib:two-axis-support}.

There is no need to recover the assigned symbols in full.  Unfold
the finite symbol-level decision tree for information-set decoding,
agreement counts, and fallback-zero outcomes.  On each fixed branch
every candidate, comparison, and output is a linear expression in
the assigned initial symbols.  For two axes these expressions are
sums of compositions of the two scalar actions
\eqref{ib:two-actions}.  Their coefficients are obtained publicly
from the fixed local matrices and the finite branch choices.

A false symbol equality is detected by a fresh uniform Hadamard
form with probability \(1/2\), or by a fresh uniform multiplier
pair with probability at least \(1/4\).  Each resulting bit of a
linear expression reduces, by additivity and scalar pullback, to
a bounded number of specified bits of the initial assigned tables.
Use Lemma~\ref{ib:inner-table-prediction} on unflagged inputs and
public zero on flagged inputs.  Independent equality trials make
the chance of missing a false equality arbitrarily small; an exact
equality fails only if one of its input-bit predictions fails.
The requested final output bit is handled in the same manner.

The order of choices prevents circular error bookkeeping.  First
fix the row rounds and the complete finite symbol-level tree.
Next choose flag and equality-detection budgets for its bounded
number of possible decisions.  This fixes the number of
initial-bit predictions needed in each tested expression.  Then
choose their smaller error budgets and their repetitions.  Given
successful flags, every later specified bit has the uniform
worst-request guarantee of the previous lemma conditional on prior
answers, using fresh splits after all scalar pullbacks.  A union
bound, or conditioning successively on the preceding correct
symbol-level decisions, bounds total simulation error by
\(\sigma\).  We do not choose the same budgets by union-bounding
over the later amplified bit tree.  After all choices that full
tree is nevertheless finite and bounded.  When a promise fails,
the same finite comparisons and fallback-zero rules still define
an output; only correctness is relinquished.
\end{proof}

\subsection{Arbitrarily accurate average correction}
\label{ib:average-correction-section}

For a point \(Y\) request its \emph{outer index} means its row
edge together with its retained point; for a line \(Y\) request
it means its edge and retained ordered line.  For \(Z\) it means
the \emph{joint pair} of edge/point indices.  A request law has
density at most \(C_{\rm req}\) within its block when its outer
index probability is at most \(C_{\rm req}\) times the corresponding
uniform probability.  Every other requested mask, form, or multiplier
may be an arbitrary public function of the request.

Call a binary word \emph{structurally canonical} if it is obtained from
the canonical recipes of Lemma~\ref{lem:inner-algebra-interface} at a
specified frame, mode, progress, and arbitrary Boolean tape, followed by
the point, line, row-extension, product, and repetition constructions
above. The tape need not satisfy the final validity conditions. Such a
word nevertheless satisfies the row-code, covariance, and multilinear
mask identities and has exactly the prescribed product tables.

\begin{proposition}[Average correction from one fixed radius]
\label{ib:average-correction}
For every fixed template there is \(\rho>0\) such that the following
holds for all sufficiently large sizes.  Let a received balanced
word \(w\) satisfy \(\dist(w,\mathcal W)\le\rho\), where
\(\mathcal W\) is structurally canonical; validity of its label
is not needed here.  For every fixed \(C_{\rm req}<\infty\) and
\(\eta>0\), a bounded-bit algorithm predicts requests with outer
density at most \(C_{\rm req}\) with average error at most \(\eta\).
It has smooth potential raw reads on every fixed hypothetical answer
branch.  The radius is independent of \(C_{\rm req}\), \(\eta\),
the row-round count, and the subsequent amplification count.
\end{proposition}

\begin{proof}
Use the constants \(\nu_1,\nu_2\) already chosen, and put
\(\zeta=\min(\zeta_*,\zeta_{**})\). Use the column and inner-table
descriptions in Section~\ref{ib:uncertainty}: rows are edges for \(Y\) and edge pairs
for \(Z\). Call an initial row position uncertain if its received inner
table is more than \(\nu_1\) from its reference table.

Each block has error fraction at most \(b_\#\rho\).  If
\(u_{\rm unc}\) is the uncertain edge fraction in a uniform
column, or pair fraction for \(Z\), then
\begin{equation}
 \E u_{\rm unc}\le\frac{b_\#\rho}{\nu_1}.
 \label{ib:uncertain-average}
\end{equation}
Call columns with \(u_{\rm unc}>\zeta\) bad.  Their fraction is
at most \(b_\#\rho/(\nu_1\zeta)\).  On every other column,
choose enough rounds in the support lemmas to leave uncertainty
fraction at most \(\delta'>0\), whose value will be fixed later.
No residual-support claim is made on bad columns.

For a \(Y\) request, independently split all its slot masks,
including the line coefficient masks, and expand.  Keep its
specified Hadamard form as the inner bit request.  For each expanded
term choose independent uniform \(h\in\mathcal H\) and replace
the row/outer pair by
\begin{equation}
 (i,g,I)\longmapsto(i,gh,h^{-1}I).
 \label{ib:random-shift}
\end{equation}
Covariance preserves the true value.  For a \(Z\) request split
both slot-mask lists and \(Q\), expand, and shift its two row/outer
indices independently by \eqref{ib:random-shift}.  Keep the
specified multipliers as its inner bit request.  Let \(T_0\) be
a fixed common upper bound on the expansion count, determined by
the template, before any accuracy choices.

For each fixed original request, every term has a uniform column:
the relevant \(\mathcal H\)-action is transitive and all its
column-mask axes have been independently randomized.  In \(Z\)
the two independent shifts give the uniform pair of outer points.
Choose the radius once and for all so that
\begin{equation}
 \frac{T_0b_\#\rho}{\nu_1\zeta}\le\frac1{64}.
 \label{ib:fixed-radius}
\end{equation}
Then, for every individual request, the chance that any term uses
a bad column is at most \(1/64\).  This conclusion does not assume
that the shifted row is independent of the shifted outer point.

For every fixed shift, \eqref{ib:random-shift} is a permutation
of the complete row/outer index; the independent pair of shifts
is a permutation of the complete \(Z\) index.  Thus averaging
over an outer-density-\(C_{\rm req}\) request law preserves that
density bound.  Expanded column masks are uniform conditional on
the indices.  The mean chance of hitting the good-column residual
support in any of the terms is consequently at most
\(T_0C_{\rm req}\delta'\).  By Markov's inequality, outside an
exceptional fraction at most \(8T_0C_{\rm req}\delta'\) of
original requests, the conditional sum of those residual-hit
probabilities is at most \(1/8\).

Bit-simulate every term's row cone by
Lemma~\ref{ib:bit-simulation}, with total simulation error at most
\(1/64\).  On any nonexceptional request, the full randomized
XOR trial therefore errs with probability at most
\[
 \frac1{64}+\frac18+\frac1{64}=\frac5{32}<\frac14.
\]
Cone-specific assignments are sufficient because the support bounds
cover every consistent assignment in each cone.  Choose
\(\delta'\) so that \(8T_0C_{\rm req}\delta'\le\eta/2\).
This determines a constant row-round count.  Repeat the entire
trial independently at the fixed request, including shifts, flags,
and all predictions, and take majority.  A bounded number of
repetitions makes its error at each nonexceptional request at most
\(\eta/2\).  Adding the exceptional fraction proves average error
at most \(\eta\) from the fixed radius \eqref{ib:fixed-radius}.
This proves an average guarantee, and makes no assertion of
arbitrarily accurate correction at every prescribed outer index.
Smoothness is proved in the next lemma.
\end{proof}

\begin{lemma}[All hypothetical read slots are smooth]
\label{ib:all-branch-density}
For the tester, the corrector of
Proposition~\ref{ib:average-correction}, and the finite linear
compositions used below, use their stated bounded outer-index request
laws, with joint pair density for \(Z\). Freeze any hypothetical sequence
of hidden answer bits.  Under the original, unconditioned public random tape
and request law, the raw-read subdistribution of every potential
slot has bounded density relative to its balanced input block.
After all accuracy counts have been fixed, the sum of these bounds
over the complete finite answer tree is a constant independent of
the size.  Defaults, wrong flags, and wrong equality answers satisfy
the same assertion.  Public numerical thresholds that only select
among these fixed schedules do not change the bound.
\end{lemma}

\begin{proof}
First enumerate every potential initial row-cone \emph{path slot},
including all candidate uses, before deduplication.  A term's
initial covariance shift is the joint permutation
\eqref{ib:random-shift}.  A fixed row-star path subsequently changes
the edge label among finitely many possibilities and multiplies
\(g\) on the left by a fixed group element.  A right-star move
from label \(i\) to \(j\), for example, uses
\(g\mapsto u_j^{-1}u_i g\).  Each such map is a permutation
within a fixed edge-label class, so costs only a fixed label factor
in a density bound.  In a two-axis cone the same reasoning applies
separately to each axis and preserves the joint pair bound.
Consequently each enumerated path slot has a fixed density bound.

Deduplication, representative selection, and omission of unused
terms only decide whether an already enumerated slot is read.
Its resulting unconditioned subdistribution is dominated by the
original slot distribution.  We never condition and renormalize
on the event that this slot was chosen as a representative or that
its flag had a specified outcome.  Summing the bounds of all
potential slots, including unexecuted ones, bounds their complete
occurrence count.  A flagged, failed, or unused zero value is public
zero and requires no placeholder query at a special address.

On a frozen symbolic branch, all information-set interpolation
matrices depend only on \(H_L,H_R\) and public graph indices.
The relevant inverses exist by their maximal-minor property.
The coefficients of all selected linear expressions, including
the compositions of \eqref{ib:two-actions}, can therefore be
computed publicly.  Hidden work is only the bounded equality and
output-bit calls described in Lemma~\ref{ib:bit-simulation}.
Their specified forms and multiplier pairs may depend on every
frozen answer, public coefficient, and preceding random choice.
Fresh splits \emph{after} these pullbacks give uniform raw inner
forms or jointly uniform multiplier pairs conditional on all that
information.  Zero coefficients can be split as well, or omitted.
Fresh uniform copies give the required copy marginal.  Wrong
answers change the selected expression, but introduce no new type
of row-cone path.

The tester's schedules were explicitly smooth before replacements.
The neighbor compositions below use finite lists of commuting
affine column permutations and changes between fixed block types,
so the same argument adds only fixed type factors.  Dropping a
public operand takes a marginal of a joint law, which also
preserves a density bound.  Finally freeze all simulation and
amplification counts, unfold the full raw bit tree, and sum these
finitely many bounds.  This tree may be large, but is constant at
the chosen accuracy.  A later threshold offset that selects its
branches changes neither its slot count nor these public address
maps.  The assertion concerns the unconditioned schedules for
all branches, not correctness on unsuccessful branches.
\end{proof}

The request statement also applies to a subdistribution within a
type: pad its missing mass by the uniform law, which increases a
density bound by at most 1, and charge the original subdistribution
by the resulting error bound.  A bounded list of type or branch
choices is handled by adding their bounds and error budgets.
This does not normalize on a rare branch.  Likewise, if a mixed
product has a public operand, a joint density
\(p(x,y)\le C/(\abs X\abs Y)\) implies
\begin{equation}
 \sum_y p(x,y)\le C/\abs X.
 \label{ib:public-marginal}
\end{equation}
The surviving \(x\) is the complete edge/point index, not only
its point.  The requested form
\(u\mapsto Q(\pi(b_1u),\pi(b_2v_{\rm public}(y)))\) may depend
on the discarded index; it does not select a new surviving index.
One may forget an actual repetition-copy index and choose a fresh
uniform copy in the corresponding one-port block.  This changes
only the fixed block-weight factor, not the original-index bound.

All public supports used here have polynomial cardinality for fixed
accuracy.  Groups, retained points and lines, representation fibers,
fields, masks, forms, and copy sets are polynomial-listable; only a
constant number of independent draws is used.  Their sampling laws
are rational.  Public indices, coefficient operations, and the
bounded Boolean answer tables can be enumerated in polynomial time.
Neither a hidden field symbol nor a hidden array is processed as one
unit-cost query.

\subsection{Neighbor bits and global port recovery}
\label{ib:neighbors-recovery}

\begin{proposition}[Local flags and encoded neighbors]
\label{ib:neighbors}
Within a sufficiently small fixed radius of a valid canonical word,
one can decide its incident-edge flags with arbitrarily high fixed
reliability, and supply its own bits and those of every existing
neighbor with arbitrarily small average error for bounded-density
requests.  The number of bit probes is bounded at each accuracy.
The complete hypothetical answer tree, including wrong decisions
and missing-edge defaults, satisfies
Lemma~\ref{ib:all-branch-density}.
\end{proposition}

\begin{proof}
Reduce \(\rho\) if needed to lie below half the distance of
Proposition~\ref{ib:tester}; a nearby valid label is then unique.
Use average correction at uniform systematic edge/point positions
to read the mode and the needed critical-bit projections of
Lemma~\ref{lem:inner-algebra-interface}.  At full progress the
all-ones contraction of the relevant singleton stack is a constant
polynomial equal to its critical Boolean read.  Use a specified
additive form taking 1 to 1; its actual raw queries are smoothed by
the corrector.  Such uniform systematic requests have bounded
outer density even though the logical read itself has a sparse
public address.

Compare the progress marker with its public empty and full tensors
by uniform separable masks and bit forms at systematic points.
On the true constant tensors a false equality has a fixed gap by
successive nonzero contractions.  Corrected probes and sufficiently
many repetitions distinguish these cases with arbitrary fixed
reliability.  Together with mode and the full-progress critical
bits, including the wrap predicate, these are exactly the finite
decisions required by the path transition rules.  Explicitly, a
pre successor increases progress unless full; at full it performs
the enabled edit or has a missing wrap edge.  A post successor
decreases progress unless empty; at empty it advances the frame.
For predecessors, full post reverses the edit or has the missing
wrap edge, other post reverses a decrease, positive-progress pre
reverses an increase, and empty pre reverses a frame advance.
An increase also distinguishes empty progress for initialization.
There is no scan of the full progress index and no proposed-tape
validity test.

Fix the true Boolean transition branch of an existing edge.  The
algebraic interface expresses every updated contraction as a
bounded sum of public terms and public scalar multiples of old
contractions at invertible \(K\)-affine pullbacks of the common
target.  Slot pulls and scalars are independent of the target point
and row.  The pullbacks commute with the rectangular action, so
they commute through the common row encoder as column permutations.
Public initialization terms have publicly computable actual
parity values.  Thus each requested point \(Y\) bit is an XOR
of boundedly many input \(Y\) requests with specified masks and
forms at bounded-density indices.  For line \(Y\), the affine
map preserves the line parameter, so the identical argument applies
to the separable coefficient contractions
\eqref{ib:line-bit}.  The derivative update uses the separately
stored child and its summed new-axis stack; it does not recompute
native packing, polynomial division, or coefficient extraction
from a requested line or parity symbol.

For a \(Z\) bit insert the two \(Y\) update sums into
\eqref{ib:product-bit} and expand by bilinearity.  Two variable
terms give an available input \(Z\) block by tag compatibility;
absorb public coefficients into multipliers or masks.  The two
row/point coordinates undergo separate affine column permutations,
with their rows fixed.  This preserves their joint density even
when the two input types coincide: their two positions are still
distinct coordinates of the product domain, never identified into
a diagonal.  A term with one public operand is an arbitrary
specified binary linear form of the other \(J\)-symbol, and
\eqref{ib:public-marginal} gives the necessary complete surviving
row/point bound.  Here the public operand is its actual encoded
\(J\)-symbol, including parity values, not an unencoded target
value.  Two public operands give a public bit.

Apply Proposition~\ref{ib:average-correction} to every needed term
with a suitably smaller accuracy.  Each output type and fixed
branch has a bounded list of transformations, including radix
pulls and zero terms.  Charge their error and density bounds by
their sum before selecting a term or branch.  Charge the finite
decision errors separately, and analyze each relevant correct
branch with its true decisions fixed before intersecting the good
events.  Wrong decisions still select from the same public lists;
missing or unused values are public zero.  Therefore the
all-hypothetical-branch smoothness assertion follows from
Lemma~\ref{ib:all-branch-density}.
\end{proof}

\begin{proposition}[Deterministic global read-port recovery]
\label{ib:global-recovery}
At a sufficiently small fixed balanced distance from a canonical
word, every order-one logical read port can be recovered exactly
in deterministic polynomial time.  In particular this applies to
the one-port template.  No constant-query claim is made for this
global decoder.
\end{proposition}

\begin{proof}
At every systematic edge/point position in the port's order-one
point block, read the entire Hadamard table including copies and
choose a nearest \(J\)-symbol, breaking ties publicly.  Whenever
the table error is less than \(1/4\), the selected symbol is true.
Markov's inequality bounds the fraction of incorrect systematic
symbols by a fixed multiple of balanced word error, accounting for
the block weight and the systematic row fraction.

Map these symbols to their represented target arguments using the
known row action.  All fibers have equal size: for a fixed target
\(y\), any full-rank \(A\) and any \(z\) determine the unique
offset \(y-Az\), and the number of choices is the same for every
row group.  Take a plurality in each fiber.  Every wrong plurality
requires at least half its fiber's symbols to be wrong, so this
loses at most a factor two.  Project outputs outside the subfield
to any fixed \(\F_q\)-estimate, which cannot spoil a correct
symbol, and take another plurality over ignored target coordinates.
After reducing the fixed radius, we have an estimate
\(\widehat f:\F_q^\ell\to\F_q\) of the native order-one
polynomial with symbol error at most a sufficiently small fixed
\(\alpha>0\).

For each desired grid point \(x\), enumerate all directions
\(z\in\F_q^\ell\).  Along the \(q\)-word
\((\widehat f(x+tz))_{t\in\F_q}\), the expected number of
errors over \(z\) is at most
\(1+(q-1)\alpha\): every nonzero \(t\) makes \(x+tz\) uniform.
Let
\[
 t_{\rm RS}=\left\lfloor\frac{q-D-1}{2}\right\rfloor.
\]
For sufficiently small fixed \(\alpha\) and large size, Markov's
inequality shows that a strict majority of directions have at most
\(t_{\rm RS}\) errors, since \(D/q\to0\).

Use the Berlekamp--Welch reconstruction method in the linear-system
presentation of Gemmell and Sudan~\cite[Appendix~A]{GemmellSudan1992}.
Specifically, uniquely decode each line by solving the linear equations
\[
 N(t)=E(t)\widehat f(x+tz)\quad(t\in\F_q),\qquad
 E\text{ monic of degree }t_{\rm RS},\quad
 \deg N\le D+t_{\rm RS}.
\]
On a good direction, the error locator, padded by any monic factor
to degree \(t_{\rm RS}\), provides a solution with
\(N=Ef_{\rm line}\).  Every solution has this equality: the
difference \(N-Ef_{\rm line}\) has degree at most
\(D+t_{\rm RS}\) and vanishes at at least
\(q-t_{\rm RS}>D+t_{\rm RS}\) correct positions.
Divide, verify degree and agreement, and use public zero on a
failure.  The strict majority of returned values at parameter zero
is exactly \(f(x)\).  Repeat for all port grid indices.
All domains, tables, directions, fields, and linear systems here
are polynomial-size, so the complete algorithm is deterministic
polynomial time.
\end{proof}

\subsection{The binary interface}

\begin{lemma}[Inner binary interface]
\label{lem:inner-binary-interface}
For the inner algebraic families of
Lemma~\ref{lem:inner-algebra-interface}, there is a deterministic
polynomial-time binary encoding with the following properties,
uniform over the circuits in each fixed padded template.
\begin{enumerate}[label=(\roman*)]
\item Its balanced blocks are the point and line \(Y\)'s and the
compatible point-product \(Z\)'s above.  Canonical valid extended
labels have constant distance, and admit the perfectly complete
bounded-bit tester and fixed-gap proximity conclusion of
Proposition~\ref{ib:tester}.
\item One radius \(\rho>0\), fixed by the template before any
request accuracy or density, permits every fixed average error
\(\eta>0\) for requests of any fixed bounded outer-index density.
For \(Z\), the requirement is a joint pair-density bound.  Other
masks, multipliers, and forms may be arbitrary request-dependent
ones.  Finite type subdistributions are charged without
renormalization.  Mixed products with a public operand may be
reduced to arbitrary linear forms using
\eqref{ib:public-marginal}.
\item In a small fixed neighborhood of a valid label, incident
flags and current or existing-neighbor encoded bits have the
bounded-query reliability and averaged accuracy in
Proposition~\ref{ib:neighbors}.  Every routine is total with
bounded decisions and public-zero defaults.
\item After all fixed accuracy counts are chosen, all hypothetical
raw read slots in the full answer tree have a finite summed
density bound as in Lemma~\ref{ib:all-branch-density}.  This
includes potential cone paths before deduplication, unsuccessful
flags and comparisons, and arbitrary pulled-back forms.  All
coefficient and address computations are public, all remaining
hidden processing uses boundedly many bits, and all random
supports and rational sampling tables are polynomial-time
constructible.
\item The common row encoder, field presentations, positive-tag
original layouts, and fixed-label formulas are identical in
corresponding actual and one-port families.  Repetition copies
may differ and are sampled afresh in the receiving family.
At a fixed distinguished source or its first successor, a
positive \(Y\) bit has Boolean degree at most three in that
port's read input, and a same-positive-tag \(Z\) bit has degree
at most six.  An order-one port has the deterministic polynomial
global recovery of Proposition~\ref{ib:global-recovery}.
\end{enumerate}
In particular, the radius for the one-port template can be fixed
before a later actual port-count bound is chosen.  The actual
template's own radius may depend on that bound through its number
of blocks.
\end{lemma}

\begin{proof}
The construction and the preceding propositions prove
(i)--(iv) and global recovery.  For (v), field and row-code choices
were made solely from the common point parameters.  The positive
charts, including per-primary verification auxiliaries and their
tagged descendants, have the literal common-index ancestry of
Lemma~\ref{lem:inner-algebra-interface}.  Columnwise linear
extension and fixed-parameter line restriction preserve that
identity.  Passing through an original index and choosing a fresh
copy therefore suffices even when balanced multiplicities differ.

At either of the two specified fixed labels, legacy, work, frame,
mode, and progress are public, and every positive persistent recipe
has Boolean degree at most three in its port input.  Public matrix
operations, division by fixed monic polynomials, coefficient
extraction, slot contraction, fixed affine substitutions, line
restriction, \(J\)-linear row extension, coordinate changes, and
additive bit forms are all \(\F_2\)-linear in these hidden-data
functions.  None raises that degree.  The one additional binary
bilinear operation in \eqref{ib:product-bit} gives degree at most
six.  Finally \eqref{ib:fixed-radius} uses only the template's
\(T_0,b_\#,\nu_1,\zeta\); the later requested density and
accuracy affect only row rounds and repetitions.  Applying it
first to the fixed one-port template proves the last assertion.
\end{proof}
